\section*{第 \ref{ch:b2:potential-wells-tunneling} 章——势阱、势垒与隧道效应}

\begin{solution}{exo:b2:potential-wells-tunneling:1}
\qty{1}{nm}: \qty{0.38}{}, \qty{1.50}{}, \qty{3.38}{eV}; $2 \to 1$: \qty{1.13}{eV}, \qty{1.1}{\micro m}.
\qty{0.1}{nm}：\qty{38}{eV}。核子：\qty{8}{MeV}。弹珠：$E_1 = \qty{5.5e-63}{J}$；
$n = 2mvL/h = 3 \times 10^{27}$.
\end{solution}

\begin{solution}{exo:b2:potential-wells-tunneling:2}
约束能 $\qty{3.76}{eV.nm^2}/L^2$：$0.10$，$0.24$，$0.42$，\qty{0.60}{eV}；总能量
$1.84$, $1.98$, $2.16$, \qty{2.34}{eV}: \qty{673}{}, \qty{628}{}, \qty{574}{}, \qty{530}{nm} ---
红、橙、黄绿、绿。
\end{solution}

\begin{solution}{exo:b2:potential-wells-tunneling:3}
$\kappa = \qty{1.09e10}{m^{-1}}$; $\eu^{-2\kappa a}$: $1.5 \times 10^{-3}$, $1.8 \times 10^{-5}$,
$3 \times 10^{-10}$，$10^{-19}$；$\times 8.8$（每 \qty{0.1}{nm}）；质子：$\kappa$ 大
$43$ 倍，\qty{0.3}{nm} 处为 $\eu^{-280}$。
\end{solution}

\begin{solution}{exo:b2:potential-wells-tunneling:4}
$k_1 = \qty{7.3e9}{m^{-1}}$，$k_2 = \qty{5.1e9}{m^{-1}}$；$R = 0.03$，$T = 0.97$。台阶
向下：$k_2 = \qty{8.9e9}{m^{-1}}$，$R = 0.01$。经典粒子在有能量时
从不折回；光波
在任何折射率变化处都会部分反射。
\end{solution}

\begin{solution}{exo:b2:potential-wells-tunneling:5}
(a)、(b) 见\cref{prop:b2:potential-wells-tunneling:infinite}。(c) $L/2$；
$0.18L$；$0$；$\pi\hbar/L$；乘积 $0.57\hbar > \hbar/2$。(d) 快速
振荡，围绕平均值 $1/L$——即经典的均匀密度。
\end{solution}

\begin{solution}{exo:b2:potential-wells-tunneling:6}
(a) 见\cref{prop:b2:potential-wells-tunneling:finite}。(b) $1$，$3$，$7$。
(c) $R = 2.6$，两个态；$\kappa = \qty{4.4e9}{m^{-1}}$，深度 \qty{0.23}{nm}。(d)
无穷多个，从下方趋于无限深势阱的情形。
\end{solution}

\begin{solution}{exo:b2:potential-wells-tunneling:7}
(a) $1 + r = t$，$\iu k(1 - r) = -\kappa t$。(b) 分子与分母
互为共轭；相位 $-2\arctan(\kappa/k)$。(c) 对实指数函数，$\psi^*\psi'$ 为
实数。(d) \qty{0.2}{nm}；与\omterm{thm:b2:wave-interfaces:snell}{全反射}的\omterm{def:b2:dispersion-wave-packets:complex}{倏逝波}相似：
存在，不携带通量，并且能够向第二种介质提供能量
——这就是\omterm{thm:b2:potential-wells-tunneling:tunnel}{隧道效应}。
\end{solution}

\begin{solution}{exo:b2:potential-wells-tunneling:8}
(a) 见\cref{thm:b2:potential-wells-tunneling:tunnel}。(b) $\kappa a = 3$：
精确值 $1/(1 + \sinh^23) = 9.9 \times 10^{-3}$，近似值 $4\eu^{-6} = 9.9 \times
10^{-3}$；$\kappa a = 1$：$0.42$ 对 $0.54$。(c) $\kappa \to \iu k_2$；$T = 1$，
条件为 $k_2a = n\pi$。(d) 两个界面的反射在
势垒宽度恰为半波长的整数倍时相互抵消——共振
透射。
\end{solution}

\begin{solution}{exo:b2:potential-wells-tunneling:9}
(a) $b = \qty{52}{fm}$; $G = 82$; $T = \eu^{-82} \approx 2 \times 10^{-36}$. (b) $v =
\qty{1.55e7}{m/s}$，每秒 $10^{21}$ 次撞击：速率 $2 \times 10^{-15}\,\unit{s^{-1}}$，
寿命 $\sim 10^7$ 年。(c) $G = 100$ 与 $69$：$\eu^{18} \approx 10^8$ 倍更长，
$\eu^{-13} \approx 10^{-6}$ 倍更短。(d) $\alpha$ 粒子束缚很紧，并
以正能量逸出；质子则不能；$^{12}$C 的电荷是其三倍，
指数也远大得多（团簇衰变确实存在，约为 $10^{-10}$）。
\end{solution}

\begin{solution}{exo:b2:potential-wells-tunneling:10}
(a) \qty{24}{GHz}; \qty{41}{ps}. (b) $\psi = (\varphi_+\eu^{-\iu E_+t/\hbar} +
\varphi_-\eu^{-\iu E_-t/\hbar})/\sqrt2$; $P = \sin^2(\pi\nu t)$. (c) $\eu^{(\sqrt2 - 1)
\kappa a} = 15$：$\kappa a = 6.5$。(d) 较重的基团与较高的势垒使
\omterm{thm:b2:potential-wells-tunneling:tunnel}{隧道效应}的周期长于数年。
\end{solution}

\begin{solution}{exo:b2:potential-wells-tunneling:11}
(a) \qty{1.09e10}{m^{-1}}；$\times 8.8$；$1\%$ 的 $I$ 对应 \qty{0.5}{pm}。(b) $6 \times 10^9$。
(c) 每个相邻原子位于 $d' = \qty{0.56}{nm}$ 处，携带的电流为针尖电流的 $\eu^{-2\kappa \times 0.06\,\text{nm}}
= 0.28$；针尖本身约携带一半。(d)
$10^{-5}\,\unit{K}$——不可能；针尖与样品同步膨胀的对称设计，
快速扫描、低膨胀材料，以及\omterm{def:b2:feedback-oscillators:loop}{反馈}。
\end{solution}

\begin{solution}{exo:b2:potential-wells-tunneling:12}
(a) \qty{0.056}{} 与 \qty{0.008}{eV}。(b) \qty{1.48}{eV}，\qty{835}{nm}。(c) 约束能
\qty{0.17}{eV}：$L = \qty{6.1}{nm}$。(d) $R = 3.6$：三个。
\end{solution}

\begin{solution}{pb:b2:potential-wells-tunneling:1}
\textbf{1.} $\kappa = \sqrt{2m\Phi}/\hbar = \qty{1.09e10}{m^{-1}}$.

\textbf{2.} $\eu^{2.2} = 8.8$；$\eu^{0.022}$：每皮米 $2.2\%$。

\textbf{3.} \qty{0.5}{pm}.

\textbf{4.} 每秒 $6 \times 10^9$ 个；相隔 \qty{0.16}{ns}，而 $\hbar/\Phi
\sim 10^{-16}\,\unit{s}$：一次一个。

\textbf{5.} 约一半；其余电流随距离衰减得如此之快，以至于
所成的像是单个原子的。

\textbf{6.} $\times 8.8$；恒定电流使针尖安全，并将
指数关系转化为线性的高度。

\textbf{7.} $\Delta T = 10^{-12}/(10^{-5} \times 0.02) = \qty{5e-6}{K}$；对称的
安装方式、扫描速度、低膨胀材料，以及能够跟随
漂移的\omterm{def:b2:feedback-oscillators:loop}{反馈}。

\textbf{8.} 梯形势垒，在收集一侧降低：有效
势垒变薄，$I$ 的增长快于线性。

\textbf{9.} $\kappa' = \qty{0.89e10}{m^{-1}}$；在 \qty{0.5}{nm} 处，指数由
$10.9$ 降至 $8.9$，$I \times 7$：针尖回缩 $\ln 7/2\kappa' \approx \qty{0.1}{nm}$
——这是一个电子效应引起的凸起，而非几何起伏。

\textbf{10.} $V(R_{\text{n}}) = 2 \times 90 \times 1.44/8 = \qty{32}{MeV}$; $b = \qty{62}{fm}$.

\textbf{11.} $v = \qty{1.4e7}{m/s}$；每秒 $9 \times 10^{20}$ 次撞击。

\textbf{12.} $G = 96$; $T = \eu^{-96} \approx 10^{-42}$.

\textbf{13.} 衰变率 $\sim 10^{-21}\,\unit{s^{-1}}$，半衰期 $\sim 10^{13}$ 年——约为
实测值的一千倍：指数为一百时，一个
粗糙的模型能落在几个数量级之内，已属不易。

\textbf{14.} $G = 79$：半衰期缩短 $\eu^{17} \approx 10^7$ 倍，约为 $10^6$ 年——
多出一个 MeV，相差七个数量级。

\textbf{15.} $G = 92$：寿命缩短 $\eu^{4.5} \approx 90$ 倍——半径
改变一飞米，相差两个数量级。

\textbf{16.} 质子并非预先具有正能量（它被
约 \qty{7}{MeV} 束缚）；$^{12}$C 的电荷为三倍，指数也
大三倍。

\textbf{17.} 是。$\alpha$ 粒子总是具有 \qty{4.2}{MeV} 的能量；其\omterm{def:b2:schrodinger-wave-functions:psi}{波函数}
仅在势垒之下有倏逝的部分，那里任何测量都不会
发现它具有负动能。

\textbf{18.} \qty{23.7}{GHz}; \qty{1.27}{cm}.

\textbf{19.} $(|{\uparrow}\rangle \pm |{\downarrow}\rangle)/\sqrt2$；
对称的那个较低——无节点，曲率较小。

\textbf{20.} $\psi = (\varphi_+\eu^{-\iu E_+t/\hbar} + \varphi_-\eu^{-\iu E_-t/\hbar})/\sqrt2$;
$P_{\downarrow} = \sin^2(\pi\nu t)$；周期 \qty{42}{ps}。

\textbf{21.} $\eu^{-\Delta/k_BT} = 0.996$：几乎相等，无\omterm{prop:b2:laser:gain}{粒子数反转}；一个
非均匀电场将分子束分选，把上能级的分子
送入谐振腔。

\textbf{22.} \qty{1.6e-23}{J}；每秒 $6 \times 10^{12}$ 个分子。

\textbf{23.} $\kappa a = 6.5$.

\textbf{24.} 它仅由分子本身决定，对每个
分子都相同，且对外界不敏感；频移来自
分子束中的\omterm{thm:b2:sound-waves:doppler}{多普勒效应}、杂散电场、碰撞，以及
谐振腔的牵引效应。

\textbf{25.} 限域使能量量子化（$E_n \propto n^2/L^2$）；倏逝
让波穿透（$T \sim \eu^{-2\kappa a}$）；指数由
$\sqrt{2m(V_0 - E)}/\hbar$ 与宽度决定——质量、高度与距离。
\end{solution}
